\documentclass{article}
\usepackage{graphicx} 
\usepackage[margin=2.5cm]{geometry}
\usepackage{natbib}
\usepackage{tikz}
\usepackage{booktabs}

\usepackage{longtable}
\usepackage{booktabs}
\usepackage{multirow}
\usepackage{array}

% Shortcuts %
\newcommand{\agents}{A}
\newcommand{\objects}{H}
\newcommand{\match}{\mu}
\newcommand{\linearorders}{\mathcal{L}}


\newtheorem{example}{Example}
\newtheorem{theorem}{Theorem}
\newtheorem{definition}{Definition}
\newtheorem{observation}{Observation}
\newtheorem{proposition}{Proposition}
\newtheorem{conjecture}{Conjecture}
\newtheorem{corollary}{Corollary}
\newtheorem{lemma}{Lemma}

%TikZ commands
\newcommand*\encircled[1]{\tikz[baseline=(char.base)]{
            \node[shape=circle,draw,inner sep=0.25pt,minimum size=1em] (char) {\vphantom{Ag}#1};}}
\newcommand*\framed[1]{\tikz[baseline=(char.base)]{
            \node[draw,inner sep=0.25pt,minimum size=1em] (char) {\vphantom{Ag}#1};}}
            

\title{Report TER}
\author{Julie Dornat}
\date{April 2026}

\begin{document}

\maketitle

\section{Abstract}

Collective decision-making problems often require allocating resources among agents whose preferences are only partially known. 
Reasoning about fairness and efficiency under such incomplete information is both practically motivated and theoretically challenging, as classical criteria can no longer be directly evaluated.

In this report, we survey existing results on fair division and matching under partial preferences, focusing on the possible and necessary approach. 
We record type of partial preferences representation and review complexity results for envy-freeness under incomplete ordinal preferences. 
Then we study popular matchings as an efficiency criterion for the house allocation problem.

Building on the characterization of \cite{abraham2007popular}, we define possible and necessary variants of popularity under partial preferences and formalize the associated computational problems. 
As a first step, we consider a simplified setting where agents have strict linear order preferences and a perfect matching is assumed to exist.
These results lay the groundwork for designing algorithms that handle popularity under incomplete information.

\section{Introduction}

Computational social choice~\citep{brandt2016handbook} studies collective decision-making problems involving multiple agents with potentially conflicting preferences. 
Its goal is to design mechanisms that aggregate individual preferences into a collective decision while satisfying desirable properties such as fairness and efficiency. 
Typical fairness criteria include maxmin and envy-freeness, which aim to capture the extent to which agents are satisfied with the outcome.

A central challenge in computational social choice lies in the elicitation and representation of agents' preferences over a set of possible alternatives. 
In many practical settings, the number of possible allocations grows combinatorially with the number of resources, making complete preference elicitation both cognitively demanding for agents and computationally challenging. 
As a result, a significant line of research focuses on developing compact and expressive preference models that mitigate these difficulties.

At the same time, collective decision-making must balance multiple, often conflicting objectives. 
As discussed in the \citep[Chapter 12]{brandt2016handbook}, fairness criteria such as envy-freeness or maximin may conflict with efficiency requirements such as Pareto efficiency or utilitarian optimality, which aim at maximizing the overall welfare and avoiding loss of resources. 
This results in seeking solution able to handle the tradeoff between fairness and efficiency.

This motivates the study of preference representations that capture only partial information, while still allowing one to reason about fairness and efficiency properties of the resulting outcomes.
In this context, partial preferences provide a natural framework to model incomplete knowledge about agents' preferences. 
As emphasized by \cite{kenig2019approximate}, incomplete preferences constitute a distinct form of uncertainty.

As discussed by \cite{lang2020collective}, this raises a fundamental question: how can one make reliable collective decisions when the available preference information is incomplete? and when can the eliciting process be stopped to deduce reliable solution?
To address this issue, the notions of possible and necessary outcomes have been introduced, allowing one to reason about decisions that hold under some or all completions of the missing information.

Computational social choice has applications in a wide range of domains, including voting, fair division, matching, multi-agent systems, and information retrieval. 
In this work, we focus on fair division and matching problems.

Fair division concerns the allocation of resources among agents in a fair manner, and applies to both divisible and indivisible goods. 
A classical example in the indivisible setting is the house allocation problem.
Such settings can also be viewed through the lens of matching theory, that provides a rich framework for modeling allocation mechanisms. 

\paragraph{Research Question}

This report investigates the following question:

\begin{quote}
How can we design fair allocation or matching mechanisms when agents' preferences are only partially known?
\end{quote}

More specifically, we explore how notions such as possible and necessary criteria can be used to reason about fairness under incomplete preferences, and how these ideas can be connected to matching frameworks.


\section{Preference Models}

Partial preferences aim to represent situations in which agents do not fully specify their preferences over all possible alternatives. 
Instead of a complete description, only partial information is available, which induces uncertainty regarding the evaluation of outcomes with respect to properties such as fairness or efficiency.

From an informational perspective, partial preferences can be seen as a compact representation of agents' preferences. 
This is closely related to the notion of representation and compilation complexity discussed in the \cite[Chapter 10, section 10.7]{brandt2016handbook}, where the goal is to encode preference information in a concise way while preserving the ability to reason about it.

During our literature review, we identified several families of partial preference models \citep[Chapter~10]{brandt2016handbook}, summarized in Appendix~\ref{tab:preference-models}. 
Among these, three models are of particular relevance to this work.

\begin{definition}[Cardinal preferences]
Agents assign numerical utilities to alternatives, capturing the intensity of their preferences.
\end{definition}

\begin{definition}[Ordinal preferences]
Agents provide a ranking over alternatives or subsets of alternatives.
\end{definition}

\begin{definition}[Pairwise comparison preferences]
Agents indicate which alternative is preferred between two options, yielding a set of binary comparisons.
\end{definition}

These models illustrate different ways of reducing the amount of information required from agents while maintaining the ability to reason about collective decisions. 
Not all these models are equally relevant for the study of fairness under incomplete information, which motivates focusing on specific classes in the remainder of this work.

In practice, the choice of a preference model has a direct impact on both the expressiveness of the representation and the computational complexity of the associated problems. 
It therefore depends on the target application and the type of reasoning required. 
Simpler models reduce the elicitation burden and facilitate computation, while richer models allow for more expressiveness at the cost of increased complexity.

A key advantage of partial preferences is that they reduce the amount of information agents need to communicate, which can improve the reliability of the elicited data while maintaining sufficient structure for decision-making.

\section{Fair Division with Partial Preferences}

Fair division is a combinatorial problem that consists in allocating a set of objects among a set of agents while satisfying certain fairness criteria. 
The nature of the agents' preferences directly impacts the definition of the problem.

While early works often assume cardinal preferences, this approach is limited in practice, as agents may not be able to provide precise numerical utilities. 
In many situations, agents can only express qualitative comparisons (e.g., preferring one object to another) or partial information (e.g., intervals of utilities).
This has led to the study of ordinal preferences, where agents provide full rankings instead. 
However, eliciting complete rankings can be impractical in settings with a large number of objects, which motivates the study of partial preferences.

In this setting, agents typically report ordinal preferences over individual goods, from which preferences over bundles must be inferred. 
A key contribution in this direction is due to \cite{brams2005efficient} and \cite{brams2003fair}, who introduce a pairwise dominance relation to compare bundles based on preferences over individual items. 
This idea is closely related to the responsive set extension, which provides a standard way to lift preferences from individual goods to bundles. 

Such extensions form the semantic foundation for defining envy between agents when only preferences over single items are available. 
Therefore, preferences over single items can be interpreted as a compact representation of preferences over bundles.
Moreover, when preferences over individual items are only partially known, this induced preferences over bundles becomes themselves incomplete, which is precisely what motivates the study of possible and necessary variants of envy-freeness.

\paragraph{Fairness Criteria under Incomplete Preferences} When preferences are incomplete, classical fairness notions such as envy-freeness can no longer be directly evaluated. 
Instead, fairness must be defined with respect to all possible completions of the agents' preferences.

\cite{bouveret2010fair} and \cite{aziz2024computational} define the notions of possible and necessary envy-freeness as follows:

\begin{definition}[Envy-Freeness]
An allocation is envy-free if no agent prefers the bundle assigned to another agent over her own bundle.
\end{definition}

\begin{definition}[Possible Envy-Freeness (PEF)]
An allocation is possibly envy-free if there exists at least one completion of the incomplete preferences under which the allocation is envy-free.
\end{definition}

\begin{definition}[Necessary Envy-Freeness (NEF)]
An allocation is necessarily envy-free if it is envy-free under all completions of the incomplete preferences.
In other words, an allocation satisfies NEF if it is envy-free for all additive valuations that are consistent with the agents' ordinal preferences over individual goods.
\end{definition}


\paragraph{Computational Aspects} The introduction of possible and necessary fairness leads to new computational questions. 
In particular, one may ask whether there exists an allocation satisfying a given fairness notion under incomplete information.

\cite{bouveret2010fair} study fair division under ordinal preferences with incomplete information, assuming that agents report only their preferences over individual goods. 
They show that:
\begin{itemize}
    \item deciding whether a possibly envy-free allocation exists is computationally simple;
    \item deciding whether a necessarily envy-free allocation exists is NP-complete in general;
    \item the problem becomes polynomial-time solvable when the number of agents is restricted to two.
\end{itemize}

These results highlight the intrinsic difficulty of reasoning under incomplete preferences, as necessary fairness requires considering all possible completions of the preference profile.
Necessarily envy-freeness is a particularly strong requirement under incomplete preferences.
It ensures that an allocation remains envy-free regardless of how the missing preference information is completed.

\cite{aziz2024computational} refine and strengthen these results by proving that the problem of deciding the existence of a necessarily envy-free allocation (ExistsNEF) remains NP-complete for any fixed number of agents greater than or equal to three and polynomial-time solvable for two agents.

\paragraph{Relaxations of Envy-Freeness} \cite{aziz2024computational} also provide several relaxations of envy-freeness.
One prominent example is envy-freeness up to one item (EF1) and its necessarily envy-free counterpart (NEF1):

\begin{definition}[EF1 (resp. NEF1)]
An allocation satisfies EF1 (resp.  NEF1) if, any envy can be eliminated by removing at most one item from the bundle of the envied agent, for every pair of agents and under some (resp. all) completions of the preferences.
\end{definition}

While these notions provide meaningful relaxations of envy-freeness, the computational complexity of determining the existence of such allocations has not been characterized in the literature. 
\cite{aziz2024computational} show that a EF1 allocation can be obtained via a round-robin procedure.

\paragraph{Conclusion} These results highlight a fundamental trade-off between criteria satisfaction and computational complexity in fair division under incomplete preferences. 
While weaker notions such as possible envy-freeness are easier to satisfy and verify, stronger notions such as necessary envy-freeness provide strong guarantees but are computationally much more demanding.

This motivates the study of intermediate notions, such as EF1 and NEF1, which are weaken fairness requirements while preserving meaningful guarantees under incomplete preferences.

In the following, we investigate how similar ideas can be adapted to matching problems, in particular through the notion of popularity under incomplete preferences.


\section{Matching and Popularity}

A particular case in matching theory, widely studied in the context of fairness, is the one-sided matching model, where only one side of the market holds preferences over the other. 
This is the case of the house allocation problem, where agents have preferences over houses but houses do not have preferences over agents, and each agent receives exactly one house. 

\paragraph{Optimality Criteria}

Matching problems admit several optimality criteria, as surveyed by \cite{cseh2017popular}.

A matching is \emph{Pareto-optimal} if no agent can be made better off without making another agent worse off. 
Introduced by \cite{gardenfors1975match}, \emph{Popular matching} offers an alternative optimality notion suited to one-sided models, satisfying a trade-off between the number of agents matched and the respect of preferences.
A matching is popular if no other matching is preferred by a more agents (see the formal definition in the Preliminaries). 

Aziz et al.~(2013) study the relationship between popularity and envy-freeness, and show that the two criteria are incompatible in general when the number of agents exceeds three. 
However, when both can be satisfied simultaneously, such a matching can be found in polynomial time. 
This tension between popularity and fairness motivates studying them independently under incomplete preferences.

\paragraph{Algorithms and Complexity}

A key challenge is that popular matchings are not guaranteed to exist. 
\cite{cseh2017popular} illustrates this with an example of the voting paradox of Condorcet.
The matchings$\mu_1$, $\mu_2$, and $\mu_3$ can form a cycle in which each is preferred by a more agents over the next, preventing any of them from being popular.

\cite{abraham2007popular} provide a algorithmic treatment of the one-sided model. 
Their algorithm either finds a popular matching of largest cardinality or certifies that none exists. 
For strict ordinal preferences, they give an $O(n+m)$ algorithm and for the general case where preference lists may contain ties, they give an $O(\sqrt{n}\, m)$ algorithm, where $n$ is the total number of agents and objects and $m$ is the total length of all preference lists.

The algorithm of \cite{abraham2007popular} relies on two key notions: the \emph{fpost} of an agent, defined as her most preferred object, and her \emph{spost}, defined as her second-most preferred object among those not ranked first by any agent. 
Computing these notions requires complete knowledge of the top of the preference ordering. 
Under partial preferences, the fpost and spost of an agent are no longer uniquely defined: they depend on how the missing preferences are completed, and may take several possible values. 
This raises algorithmic challenges.
In this sense, the characterization of \cite{abraham2007popular} serves as a structural baseline for redesigning algorithms under incomplete information.

Furthermore, the extension of \cite{abraham2007popular} to preference lists with ties provides a useful technical analogy for our setting.
Partial orders and preference orderings with ties share the property of leaving certain pairwise comparisons unresolved.
Whether the algorithmic techniques developed for ties may serve as a starting point for handling incompleteness remains an open direction.

\paragraph{Towards Partial Preferences}

All of the results above assume that agents' preferences are fully known. 
When preferences are only partially specified, it is no longer possible to directly evaluate whether a matching is popular. 
Following the approach developed for envy-freeness in the previous section, one can define possible and necessary variants of popularity with respect to the set of all completions of the partial preference profile. 
This gives rise to four computational problems which are formally defined in the Preliminaries and form the focus of this work.

\section{Preliminaries}

A House allocation instance $I$ with partial preferences is given by $<\agents, \objects, P>$ where $\agents=\{a_1, ..., a_n\}$ is a set of $n$ applicants, $\objects = \{h_1, ..., h_n\}$ is a set of $n$ houses, and $P=\{P_1, ..., P_n\}$ is the partial preference profile, i.e. the lists of incomplete preferences for all applicants. 

Each applicant $a_i$ expresses their incomplete preferences over the set of houses $\objects$ by giving a strict partial order $P_i$ over $\objects$. 
More precisely, $h_l P_i h_k$ denotes the pairwise comparison of the houses $h_l$ and $h_k$ for applicant $a_i$, yielding that applicant $a_i$ prefers the house $h_l$ to $h_k$.
The preferences are transitive, i.e if an applicant $a_i$ prefers a house $h_l$ to a house $h_k$ and $h_k$ is preferred to a house $h_z$, then she prefers $h_l$ to $h_z$.

We assume that each applicant $a_i \in \agents$ has a complete preference linear order $\succ_i$ over $\objects$, which is unobserved due to the elicitation process.
The complete preference profile $\succ$ is the collection of all applicants' preferences, i.e., $\succ=(\succ_1,\dots,\succ_n)$.

A linear order $\succ_i$ is a linear extension of $P_i$ if it is compatible with $P_i$, i.e., for all houses $h_l,h_k\in \objects$ if $h_l P_i h_k$ then $h_l \succ_i h_k$. 
A preference profile $\succ$ is a possible completion of a partial preference profile $P$ if $\succ_i$ is a linear extension of $P_i$, for every $a_i\in \agents$.

A matching is a function that assigns houses to applicants. 
A perfect matching $\match:\agents\to \objects$ is a bijection assigning to every applicant $a_i\in \agents$ a distinct house in $\objects$ denoted by $\match(a_i)$.

A matching $\match$ is \emph{Pareto-optimal} if there is no other matching that is at least as good for all applicants and strictly better for at least one applicant.
We are interested in a refinement of Pareto-optimality called \emph{popularity}: a matching is popular if there is no other matching that is more popular, i.e., that is preferred by more applicants.

\begin{definition}[Popularity]
A matching $\match$ is popular if $|\{i\in \agents: \match(i)\succ_i \match'(i)\}|\geq |\{i\in \agents: \match'(i)\succ_i \match(i)\}|$, for every matching $\match'$.
\end{definition}

A popular matching is trivially Pareto-optimal but is not guaranteed to exist.
The following characterization of popularity due to \citet{abraham2007popular} allows to decide in polynomial time whether a popular matching exists.

\begin{lemma}[\citet{abraham2007popular}]\label{lem:charact-pop}
A matching is popular iff every applicant is assigned either her most preferred item or her most preferred item among those that are not ranked first by any applicant.
\end{lemma}

The definition of popularity can then be adapted to partial preferences by considering possible or necessary outcomes with respect to the possible preference completions.

\begin{definition}[Possible popularity]
A matching is possibly popular if it is popular in at least one completion of the partial preference profile.
\end{definition}

\begin{definition}[Necessary popularity]
A matching is necessarily popular if it is popular in all completions of the partial preference profile. 
\end{definition}

We can then investigate the following computational problems:
\smallbreak
\begin{tabular}{ll}
\toprule
\multicolumn{2}{c}{\textsc{ExistPossibPop}/\textsc{ExistNecessPop}} \\
\midrule
Instance: & House allocation instance with incomplete preferences $(N,P)$ \\
Question: & Does there exist a possibly/necessarily popular matching? \\
\bottomrule
\end{tabular}

\smallbreak
\begin{tabular}{ll}
\toprule
\multicolumn{2}{c}{\textsc{VerifPossibPop}/\textsc{VerifNecessPop}} \\
\midrule
Instance: & House allocation instance with incomplete preferences $(N,P)$, matching $\match$ \\
Question: & Is $\match$ possibly/necessarily popular? \\
\bottomrule
\end{tabular}


The first problem investigates the existence of a matching that is popular under some or all completions of the partial preference profile.
The second problem verifies whether a given matching is popular. It induces checking if there exists a completion that makes the matching popular.

\section{Conclusion}


This report has investigated the problem of designing fair and efficient allocation mechanisms when agents' preferences are only partially known. 
We have surveyed two main lines of research that provide the theoretical foundation for this work.

In the context of fair division, \cite{bouveret2010fair} and \cite{aziz2024computational} have studied possible and necessary envy-freeness under ordinal preferences. 
Their results reveal an asymmetric: while possible envy-freeness is easy to decide, necessary envy-freeness is NP-complete for three or more agents, and polynomial-time solvable only for two agents. 
This pattern suggests that reasoning under all completions of a partial preference profile is intrinsically harder than reasoning under some completion.

In the context of matching, \cite{abraham2007popular} provide a complete algorithmic treatment of popular matchings under complete preferences, including an $O(n+m)$ algorithm for strict preferences and an $O(\sqrt{n}\,m)$ algorithm for preferences with ties. 
Their characterization via fpost and spost serves as a structural baseline for our setting.

Building on these two lines of work, we have defined possible and necessary variants of popularity for the house allocation problem under partial preferences, and formalized the four computational problems \textsc{ExistPossibPop}, \textsc{ExistNecessPop}, \textsc{VerifPossibPop}, and \textsc{VerifNecessPop}. 
These problems remain open, and constitute the natural next step of this research. 
In particular, the following questions are of interest:
\begin{itemize}
    \item Can the problem be solved by properties over popularity?
    \item What is the computational complexity of \textsc{ExistPossibPop}, \textsc{ExistNecessPop}, \textsc{VerifPossibPop}, and \textsc{VerifNecessPop}?
    \item Can the characterization of \cite{abraham2007popular} be used to establish the proof of correctness and coverage of all cases of algorithms for the possibly and necessarily popular matching problems?
    \item Does the asymmetry observed for envy-freeness --- possible easy, necessary hard --- also hold for popularity?
\end{itemize}

\bibliographystyle{plainnat}
\bibliography{ref}

\appendix
\section{Taxonomy of Partial Preference Models}


\begin{longtable}{p{2.5cm} p{3cm} p{8.5cm}}
\toprule
\textbf{Category} & \textbf{Model} & \textbf{Description} \\
\midrule
\endfirsthead

\toprule
\textbf{Category} & \textbf{Model} & \textbf{Description} \\
\midrule
\endhead

\bottomrule
\multicolumn{3}{r}{\small\textit{Continued on next page}}
\endfoot

\bottomrule
\endlastfoot

\multirow{2}{2.5cm}{\textbf{(1) Preference Representation}}
& Ordinal & Agents provide a ranking over alternatives or 
subsets of alternatives. \\
& Cardinal & Agents assign numerical utilities to 
alternatives, capturing the intensity of preferences. \\

\midrule

\multirow{10}{2.5cm}{\textbf{(2) Partial Elicitation}}
& Truncated: Top-$k$ & Only the $k$ most preferred 
alternatives are reported. \\
& Truncated: Bottom-$k$ & Only the $k$ least preferred 
alternatives are reported, leaving the rest unranked. 
The most preferred alternative may remain unknown. \\
& Truncated: Set ranking & Agents provide a ranking of 
a subset of alternatives. \\
& Truncated: Cardinal intervals & Agents provide ranges 
of utility values instead of a single utility. \\
& Pairwise comparisons & Agents indicate which alternative 
is preferred between two options. Central in majority-based 
and Condorcet-based reasoning (e.g., Condorcet winner, 
Copeland's rule). \\
& Positional queries & Agent provides the alternative she 
ranks at a given position, leading to a scoring vector. 
\cite{kenig2019approximate} analyze the complexity of 
estimating election outcomes under uncertainty using 
such vectors. \\
& Approval & Agents classify alternatives into approved 
and non-approved sets. \\
& Veto & Agents reject a subset of alternatives. \\
& Acceptability-based & Agents specify which alternatives 
are acceptable. Not all objects are allocable to a given 
agent. Amounts to defining preferences by excluding 
vetoed alternatives. \\
& Next-best queries & Agents iteratively reveal their 
most preferred alternative among remaining ones 
(top-$k$ queries when stopped at rank $k$). \\
& Split queries & Agents partition alternatives into 
ordered groups \citep[Chapter~10]{brandt2016handbook}. 
The $k$-approval is a special case with two blocks 
of size $k$ and $m-k$. \\
& Choice queries & Agents select their most preferred 
alternative from a subset. \\

\midrule

\multirow{3}{2.5cm}{\textbf{(3) Structured Domains}}
& Single-peaked & Agents have a most preferred option; 
preferences decrease away from this peak. Assumes an 
axis on which alternatives are positioned independently 
of agents. \\
& Euclidean & Preferences derived from distances in a 
geometric space parameterized by alternative 
characteristics. \cite{bogomolnaia2007euclidean} define 
Euclidean preferences as a point in a $d$-dimensional 
space. \\
& Attribute-based & Agents have different preference 
profiles depending on the attribute considered 
(politics, economics, etc). \\

\midrule

\textbf{(4) Probabilistic}
& Probabilistic models & Preferences are drawn from a 
probability distribution (e.g., impartial culture or 
Polya-Eggenberger urn models), allowing one to reason 
under uncertainty. \\

\caption{Taxonomy of partial preference models.}
\label{tab:preference-models}
\end{longtable}


\end{document}

